\documentclass[10pt,a4paper]{article}
\usepackage[margin=20mm]{geometry}
\usepackage{amsmath,booktabs,hyperref}
\title{Execution-Aware Skill Adaptation for Manipulation:\\Benchmark Pivot and Preliminary Baselines}
\author{ARES Research Working Draft}
\date{27 October 2026 -- Not submission-ready}
\begin{document}
\maketitle
\begin{abstract}
We investigate robot manipulation through reusable feedback skills and learned
corrections. This report establishes an executable public-benchmark track and a
separate JAKA application track, rather than claiming a new algorithm before its
mechanism has been validated. Native ManiSkill stacking experiments demonstrate
that a structured feedback controller is already strong: it completes 58 of 64
development-validation cases. An ordinary residual PPO pilot completes 61 cases;
forcing the residual to zero recovers the exact feedback outcomes. The learned
controller fixes five failures and introduces two. This is preliminary evidence
of physical action influence, not evidence of statistical superiority or novelty.
\end{abstract}

\section{Problem and scope}
The intended problem is how reusable manipulation skills should adjust legal
parameters or recover after observing actual execution outcomes. A local skill
may grasp, transport, place or release an object, but reaching an intermediate
target does not ensure stable completion. We retain ordinary motion planning and
feedback control as strong components. Reinforcement learning is evaluated by
physical task success and execution cost, not merely prediction loss.

The public track preserves the default robot, dynamics, reward and success
predicate of ManiSkill tasks. It runs in SAPIEN and is not labelled an Isaac Lab
experiment. The application track uses JAKA and EG2 assets in Isaac Sim 5.1 and
Isaac Lab 2.3.0. Replacing Panda with JAKA produces an adapted task, not an
official benchmark score. Neither track demonstrates real-robot deployment.

\section{Relation to prior work}
STAP studies skill values and planning; RAPS uses parameterized action primitives;
ViSkill and Sequential Dexterity study downstream feasibility and chaining;
RecoveryChaining combines nominal skills with learned recovery. Residual RL is
also established. Consequently, parameterization, skill switching, downstream
value and residual addition are not claimed as novel here. The proposed research
increment remains a hypothesis: exploit current execution outcomes to select
legal skill parameters or bounded recovery actions under a shared budget, with
independent causal ablations. A future manuscript must document the exact
difference from these methods and cite their original papers after the method
is fixed; this document is not a completed literature review.

\section{Implemented baseline}
Let $s_t$ contain the native simulated robot and object state. A feedback
controller has phase $z_t$ and produces $u_0(s_t,z_t)$. Both matched PPO variants
observe $o_t=[s_t,\operatorname{onehot}(z_t),u_0(s_t,z_t)]$. The direct baseline
executes $u_t=\operatorname{clip}(\pi_\theta(o_t),-1,1)$; the residual baseline
executes
\[
 u_t=\operatorname{clip}\bigl(u_0(s_t,z_t)+0.15\pi_\theta(o_t),-1,1\bigr).
\]
The native Cartesian delta-position and gripper controller converts this command
to physical actuation. No object pose is overwritten during an episode. Phase
updates use current geometry and elapsed dwell, not future success labels. PPO,
GAE and the upstream network are standard. No new auxiliary head is claimed.

The native state is privileged relative to an unimplemented vision system; it
is explicitly a state-based simulation experiment. Controller phase and nominal
command are observable internal controller states. The actor does not receive
future suffix execution outcomes. The environment reward and 50-step StackCube
horizon remain unchanged.

\section{Experimental protocol}
ManiSkill is pinned to commit
\texttt{62ff3a5896b4d5b4cf0ac4c8d79afe600c9404a3}. Training uses 64 parallel
environments, 50-step rollout windows, eight minibatches and the upstream PPO
implementation. Matched Cartesian and residual pilots use 256,000 online
interactions each. A separate upstream joint-delta PPO reference uses 1,024,000
interactions. Different action interfaces prevent treating that reference as a
fair mechanism ablation.

Development validation uses 16 batches of four environments with batch seeds
$27101+4k$, $k=0,\ldots,15$. Every policy is evaluated on these same cases with
deterministic actions. These are development cases, not a sealed final test.
Per-step action, clipping, phase, state, reward and native task predicates are
retained, together with checkpoint and source hashes. Success is the native final
predicate; success-ever is reported separately in raw artifacts.

\section{Preliminary results}
\begin{center}
\begin{tabular}{lrr}
\toprule
Controller & Training interactions & Final success \\
\midrule
Structured Cartesian feedback & 0 & 58/64 \\
Direct Cartesian PPO, seed 31 & 256,000 & 0/64 \\
Ordinary residual PPO, seed 31 & 256,000 & 61/64 \\
Ordinary residual PPO, seed 37 & 256,000 & 62/64 \\
Ordinary residual PPO, seed 43 & 256,000 & 60/64 \\
Same residual checkpoint, forced zero & same checkpoint & 58/64 \\
Upstream joint-delta PPO, seed 31 & 1,024,000 & 0/64 \\
\bottomrule
\end{tabular}
\end{center}
The residual changes executed Cartesian commands and physical outcomes. It
recovers five feedback failures, loses two feedback successes, shares 56
successes and one failure. The exact paired binomial/McNemar test on the seven
discordant cases has two-sided $p=0.453125$. Thus these cases do not establish
statistical superiority. The mean norm of the normalized executed-minus-nominal
XYZ command is approximately 0.148; it is exactly zero in the zero-residual
ablation. Fixed seeds 37 and 43 complete 62 and 60 cases respectively. The
three-seed mean is 95.3125\%, with sample standard deviation 1.5625 percentage
points, versus feedback 90.625\%. These repeat the same 64 development cases;
they are not 192 independent scenes. Seven final checkpoints and 2,560,000
training interactions were completed across all pilots, excluding separately
logged evaluation/scout steps. No final checkpoint was selected by its score.

\section{Task diagnostics and application boundary}
A GPU planner-observation defect made the native scripted Stack planner use a
stale CPU-component OBB instead of the current actor pose. Correcting only this
observation adapter changed four paired diagnostic cases from failure to success.
However, they required 88--103 steps, beyond the native 50-step budget. They are
not official within-budget successes. Corrected PegInsertionSide planning also
completed four cases, predominantly beyond its 100-step horizon. These results
identify solvability and budget issues, not algorithmic novelty.

Isaac Lab Lift, Stack and Drawer were physically run with one and four
environments. Lift and Drawer configurations lack a uniform success termination;
measured height and drawer displacement must not be presented as an invented
official success rate. JAKA/EG2 drove for 720 physical steps, but independent
joint-limit and mimic checks failed. It remains an engineering smoke, not grasp
or transfer success. Historical P1/P2 and field evidence remain separate.

\section{Next discriminating study and limitations}
The immediate study compares ordinary residual control with execution-result-
conditioned legal parameter decisions, keeping observation, action constraints,
training budget and physical conditions matched. Necessary ablations remove
execution-result conditioning and force learned corrections to zero. Extra
recovery/suffix interactions count toward the budget. A second contact family
and frozen geometry/dynamics splits are required before claiming generalization.
More independent training seeds and stronger feedback tuning are required before
claiming a repeatable advantage. Current results support baseline viability,
not a publishable new method. Asset licensing and JAKA contact qualification
also remain open; restricted field data and simulation assets are not included
in this public package.
\end{document}
